%% =====================================================================
%%  Chapter 4: Proposed Method
%%  Your own contribution. Keep it separate from the existing knowledge in the
%%  Preliminaries chapter, and describe it in enough detail that a third party
%%  could implement and reproduce it.
%%  The steps and the equation below are examples. Replace them with your own.
%% =====================================================================
\chapter{Proposed Method}
\label{chap:method}

This chapter proposes a method for the YY problem using XX.
\begin{writingnote}
Start with the overall picture of the method (what it takes as input, what it outputs and which steps it consists of),
and only then describe each step in detail; in other words, show the reader the wall before the bricks~\cite{glasmandeal2010science}.
For each step, state whether you use an existing method as it is, modify an existing method or propose something new.
Cite the existing methods, and place each citation so that your own ideas are not credited to others, and vice versa.
\end{writingnote}

\section{Overview}
\label{sec:overview}

%% Showing the overall picture of the method in a figure first helps the reader.
%% Simple diagrams can be drawn directly in LaTeX with TikZ (the tikzpicture environment below).
%% For complex diagrams, use a tool such as draw.io or PowerPoint, save the result as PDF,
%% and include it with \includegraphics.
\Cref{fig:flow} shows the flow of the proposed method.
The method consists of three steps: preprocessing, construction and improvement,
and the improvement step is repeated until a termination condition is met.

\begin{figure}[tbp]
  \centering
  \begin{tikzpicture}[
      block/.style={draw, thick, fill=black!3, rounded corners, font=\small,
                    minimum width=36mm, minimum height=8mm, align=center},
      arrow/.style={-{Stealth[length=2.5mm]}, thick},
      node distance=5mm,
    ]
    \node[block] (input) {Instance};
    \node[block, below=of input] (pre) {Preprocessing};
    \node[block, below=of pre] (build) {Construction};
    \node[block, below=of build] (improve) {Improvement};
    \node[block, below=of improve] (output) {Solution};
    \draw[arrow] (input) -- (pre);
    \draw[arrow] (pre) -- (build);
    \draw[arrow] (build) -- (improve);
    \draw[arrow] (improve) -- (output);
    \draw[arrow] (improve.east) -- ++(8mm, 0) |- node[pos=0.25, right, font=\small, align=left]
      {repeat until the\\condition is met} (build.east);
  \end{tikzpicture}
  \caption{Flow of the proposed method}
  \label{fig:flow}
\end{figure}

\section{Details of Each Step}
\label{sec:procedure}

\begin{writingnote}
Give enough detail for a third party to implement the method and reproduce the results:
the inputs and outputs, the operations in each step, the parameters and the termination condition.
Justify each design choice: explain why you chose this method or evaluation function and why you rejected the alternatives.
In this chapter, also state the assumptions of the method and the cases in which it may not work well,
rather than mentioning them for the first time in the conclusion.
Leave the performance of the method and the experimental results to the Computational Experiments chapter.
\end{writingnote}

%% The description environment is handy for explaining parallel items such as steps or symbols.
The steps are as follows.
\begin{description}
  \item[Preprocessing] reduces the size of the instance by \ldots.
  \item[Construction] builds an initial solution by extending the greedy algorithm (\cref{alg:greedy}) so that \ldots.
  \item[Improvement] improves the solution by \ldots.
\end{description}

%% Number the key equations of your method and define their symbols right after them.
The improvement step maximizes the evaluation function $f$ in \cref{eq:penalty}:
\begin{equation}
  f(\mathbf{x}) = \sum_{i=1}^{n} v_i x_i - \lambda \max\Bigl\{ 0,\ \sum_{i=1}^{n} w_i x_i - W \Bigr\},
  \label{eq:penalty}
\end{equation}
where $\mathbf{x} = (x_1, \dots, x_n)$ is the vector representing a solution
and $\lambda > 0$ is the weight of the penalty for exceeding the capacity.
